%%
%% TO DO:
%%
%%  . paste in the example program?
%%
%%  . write variant of the example program with ktjet 
%%
%%  . maybe overload the exclusive routines by allowing access
%%    interfaces such as exclusive_jets(jet, dcut) to find the
%%    subjets of jet.
%% 
%%
%%  

\documentclass[12pt]{article}

\usepackage{a4wide}
\usepackage{xspace}
\usepackage{amsmath}
\usepackage{amssymb}
\usepackage{url}

\newcommand{\fastjet}{\texttt{FastJet}\xspace}
\newcommand{\ktjet}{\texttt{KtJet}\xspace}
\newcommand{\clhep}{\texttt{CLHEP}\xspace}
\newcommand{\ttt}[1]{\texttt{#1}}

\title{FastJet 1.0 user manual}
\author{Matteo Cacciari and Gavin P. Salam\\
  LPTHE, Univsersit\'e Pierre et Marie Curie -- Paris 6, \\Univsersit\'e
  Denis Diderot -- Paris 7, CNRS UMR 7589
}
\date{}

\begin{document}
\maketitle

%----------------------------------------------------------------------
\section{Introduction}

This note documents the \ttt{FastJet} package which implements
efficient geometrically-based algorithms \cite{fastjet} for the
longitudinally invariant $k_t$ jet-finder \cite{ktexcl,ktincl}.

This release has two aims: (a) to enable people with an immediate need
for speed in $k_t$ jet-finding to have immediate access to fast
algorithms and (b) to provide access, for those who are interested, to
the implementation so that they may study it in more detail than was
exposed in \cite{fastjet}. We anticipate that the \fastjet code will
also be made available at a later date with a more developed interface
as part of the planned multi-purpose \texttt{HepJet} jet-finding
package \cite{HepJet}.

There are many variants of the $k_t$ jet-finder. \fastjet implements
the so-called $\Delta R$ distance measure with $E$-scheme
recombination, as recommended in \cite{RunII-jet-physics}. It has both
exclusive \cite{ktexcl} and inclusive modes \cite{ktincl}, the latter
currently being in more widespread use, though the exclusive mode has
physical advantages in certain cases.  We provide access to both
modes.

The definition of the inclusive $k_t$ jet finder that is coded is as
follows:
\begin{itemize}
\item[1.] For each pair of particles $i$, $j$ work out the $k_t$
  distance 
  \begin{equation}
    \label{eq:dij}
    d_{ij} = \min(k_{ti}^2,{k_{tj}^2}) \Delta R_{ij}^2 / R^2
  \end{equation}
  with $\Delta R_{ij}^2 = (\eta_i-\eta_j)^2 + (\phi_i-\phi_j)^2$,
  where $k_{ti}$, $\eta_i$ and $\phi_i$ are the transverse momentum,
  rapidity and azimuth of particle $i$ and $R$ is a jet-radius
  parameter usually taken of order $1$; for each parton $i$ also work
  out the beam distance $d_{iB} = k_{ti}^2$.
\item[2.] Find the minimum $d_{\min}$ of all the $d_{ij},d_{iB}$. If
  $d_{\min}$ is a $d_{ij}$ merge particles $i$ and $j$ into a single
  particle, summing their four-momenta (this is $E$-scheme
  recombination); if it is a $d_{iB}$ then declare particle $i$ to be
  a final jet and remove it from the list.
\item[3.] Repeat from step 1 until no particles are left.
\end{itemize}
The exclusive jet algorithm is similar except that (a) when a $d_{iB}$
is the smallest value, that particle is considered to become part of
the beam jet (i.e.\ is discarded) and (b) clustering is stopped when
all $d_{ij}$ and $d_{iB}$ are above some $d_{cut}$. Generally in the
exclusive mode $R$ is set to $1$.

%----------------------------------------------------------------------
\section{Interface}

The \fastjet code is written in C++. From the point of view of simple
usage its interface has a number of similarities to that of
\ttt{KtJet} \cite{KtJet}, fundamental differences being highlighted
below.

As with \ktjet the user is exposed to two main classes:
\begin{verbatim}
  class FjPseudoJet;
  class FjClusterSequence;
\end{verbatim}
\ttt{FjPseudoJet} provides a jet object with a four-momentum and some
internal indices to situate it in the context of a jet-clustering
sequence. \ttt{FjClusterSequence} is the class that carries out
jet-clustering and provides access to the final jets.


%......................................................................
\subsection{FjPseudoJet}
\label{sec:FjPseudoJet}


All jets, as well as input particles to the clustering (optionally)
are \ttt{FjPseudoJet} objects.  They can be created using one of the
following initializers
\begin{verbatim}
  FjPseudoJet (double px, double py, double pz, double  E);
  template<class T> FjPseudoJet (const T & some_lorentz_vector);
\end{verbatim}
where the second form allows the initialisation to be obtained from
any class \ttt{T} that allows subscripting to return the components of
the momentum (running from $0\ldots3$ in the order $p_x,p_y,p_z,E$),
for example the \ttt{CLHEP} \ttt{HepLorentzVector} class.\footnote{
  \ttt{FjPseudoJet} is the analogue of \ktjet's \ttt{KtLorentzVector}.
  A significant difference is that it is not derived from
  \ttt{HepLorentzVector} (so as to allow compilation even without
  \clhep).}
%
It has the following member functions for accessing the components
\begin{verbatim}
  double E()     const ; // returns the energy component
  double e()     const ; // returns the energy component
  double px()    const ; // returns the x momentum component
  double py()    const ; // returns the y momentum component
  double pz()    const ; // returns the z momentum component
  double phi()   const ; // returns the azimuthal angle in the range 0--2pi
  double rap()   const ; // returns the rapidity
  double rapidity() const ; // returns the rapidity
  double kt2()   const ; // returns the squared transverse momentum
  double perp2() const ; // returns the squared transverse momentum
  double perp()  const ; // returns the transverse momentum
  double operator[] (int i) const; // returns component i

  /// return a valarray containing the four-momentum (components 0--2
  /// are 3-momentum, component 3 is energy).
  valarray<double> four_mom() const;
\end{verbatim}
It also allows the user to set and access an index, in case the user
wishes to keep track of the identity of a \ttt{FjPseudoJet} object
\begin{verbatim}
  /// set the user_index, intended to allow the user to label the object
  void set_user_index(const int index);

  /// return the user_index
  int user_index() const ;
\end{verbatim}

We also provide routines for taking an unsorted vector of
\ttt{FjPseudoJet}s and returning a sorted vector,
\begin{verbatim}
  /// return a vector of jets sorted into decreasing transverse momentum
  vector<FjPseudoJet> sorted_by_pt(const vector<FjPseudoJet> & jets);
  
  /// return a vector of jets sorted into increasing rapidity
  vector<FjPseudoJet> sorted_by_rapidity(const vector<FjPseudoJet> & jets);
  
  /// return a vector of jets sorted into decreasing energy
  vector<FjPseudoJet> sorted_by_E(const vector<FjPseudoJet> & jets);
\end{verbatim}
This will typically be used on the jets returned by
\ttt{FjClusterSequence}.

%......................................................................
\subsection{FjClusterSequence}
\label{sec:FjClusterSequence}

To run the jet clustering, create a \texttt{FjClusterSequence}
object,\footnote{The analogue of \ktjet's \ttt{KtEvent}.}  through the
following initalizer
\begin{verbatim}
  template<class T> FjClusterSequence (const vector<T> & input_particles,
                   const double & R = 1.0,
                   const FjStrategy & strategy = Best);
\end{verbatim}
where \ttt{input\_particles} is the vector of initial particles of any
type (\ttt{FjPseudoJet}, \ttt{HepLorentzVector}, etc.) that can be
used to initialise a \ttt{FjPseudoJet}; \texttt{R}, which may be
omitted (and then takes value $1$) sets the value of $R$ in
eq.(\ref{eq:dij}), while \texttt{strategy}, which may also be omitted
selects the algorithmic strategy to use while clustering and is an
\texttt{enum} of type \texttt{FjStrategy} with the following four
potentially interesting values:
\begin{center}
\begin{tabular}{l|l}\hline
  \texttt{N2Plain} & a plain $N^2$ algorithm (fastest for $N \lesssim 50$)\\
  \texttt{N2Tiled} & a tiled $N^2$ algorithm (fastest for $50 \lesssim
  N \lesssim 10^4$)\\ 
  \texttt{NlnN} & the Voronoi-based $N\ln N$ algorithm (fastest for $
  N \gtrsim 10^4$)\\ 
  \texttt{Best} & automatic selection of the best of these based on $N$\\\hline 
\end{tabular}
\end{center}
If \texttt{strategy} is omitted then the \texttt{Best} option is used.

If the user wishes to access inclusive jets, the following member
function should be used
\begin{verbatim}
  /// return a vector of all jets (in the sense of the inclusive
  /// algorithm) with pt >= ptmin. 
  vector<FjPseudoJet> inclusive_jets (const double & ptmin = 0.0) const;
\end{verbatim}
where \texttt{ptmin} may be omitted (then implicitly taking value
$0$).

There are two ways of accessing exclusive jets,\footnote{In contrast
  to \ktjet the class constructor is the same for the inclusive and
  exclusive cases. This choice has been made because the clustering
  sequence is identical in the two cases.} one where one specifies
$d_{cut}$, the other where one specifies that the clustering is taken
to be stopped once it reaches the specified number of jets.
\begin{verbatim}
  /// return a vector of all jets (in the sense of the exclusive
  /// algorithm) that would be obtained when running the algorithm
  /// with the given dcut.
  vector<FjPseudoJet> exclusive_jets (const double & dcut) const;

  /// return a vector of all jets when the event is clustered (in the
  /// exclusive sense) to exactly njets.
  vector<FjPseudoJet> exclusive_jets (const int & njets) const;
\end{verbatim}
Note that these two member functions have the same name, and the
compiler selects the correct one based on the type of the arguments,
as is standard in C++.
%
The \ttt{FjPseudoJet} vectors returned by the above routines can all be
sorted with the routines described at the end of
section~\ref{sec:FjPseudoJet}.

The user may also wish just to obtain information about the number of
jets in the exclusive algorithm:
\begin{verbatim}
  /// return the number of jets (in the sense of the exclusive
  /// algorithm) that would be obtained when running the algorithm
  /// with the given dcut.
  int n_exclusive_jets (const double & dcut) const;
\end{verbatim}
Another common query is to establish the $d_{\min}$ value associated
with merging from $n+1$ to $n$ jets. Two member functions are
available for determining this:
\begin{verbatim}
  /// return the dmin corresponding to the recombination that went from
  /// n+1 to n jets (sometimes known as d_{n n+1}).
  double exclusive_dmerge (const int & njets) const;

  /// return the maximum of the dmin encountered during all recombinations 
  /// up to the one that led to an n-jet final state; identical to
  /// exclusive_dmerge, except in cases where the dmin do not increase
  /// monotonically.
  double exclusive_dmerge_max (const int & njets) const;
\end{verbatim}
The first returns the $d_{\min}$  in going from $n+1$ to $n$ jets. 
%
Occasionally however the $d_{\min}$ value does not increase
monotonically during successive mergings and using a $d_{cut}$ smaller
than the return value from \ttt{exclusive\_dmerge} does not guarantee
an event with more than \ttt{njets} jets.
%
For this reason the second function \ttt{exclusive\_dmerge\_max} is
provided --- using a $d_{cut}$ below its return value is guaranteed to
provide a final state with more than $n$ jets, while using a larger
value will return a final state with $n$ or fewer jets.


Finally the user may obtain the constituent particles of a given jet,
via
\begin{verbatim}
  /// return a vector of the particles that make up jet
  vector<FjPseudoJet> constituents (const FjPseudoJet & jet);
\end{verbatim}
Note that this is a member function of the \ttt{FjClusterSequence} and
not of \ttt{FjPseudoJet}, because jets are only meaningful when
referred to a given \ttt{FjClusterSequence}.\footnote{This is a
  further respect in which the interface differs from that of
  \ttt{KtJet}.}
%
If the user wishes to identify the constituents with the original
particles provided to \ttt{FjClusterSequence} she or he should have set
a unique index for each of the original particles with the
\ttt{FjPseudoJet::set\_user\_index} function.

\paragraph{Subjet analysis.} Currently no dedicated subjet
analysis is provided. It is however trivial (though non-optimal) to
obtain a subjet analysis as follows:
\begin{verbatim}
  // run clustering on vector<FjPseudoJet> all_particles
  FjClusterSequence clust_seq(all_particles);

  // get the exclusive jets with the event viewed as two hard jets
  vector<FjPseudoJet> jets = clust_seq.exclusive_jets(2);

  // get the constituents of the first of these jets
  vector<FjPseudoJet> jet0_constituents = clust_seq.constituents(jets[0]);

  // now rerun the jet clustering on jet0_constituents
  FjClusterSequence jet0_seq(jet0_constituents);

  // run whatever subjet analysis is of interest here on jet0_seq.
\end{verbatim}



%======================================================================
\section{Example program}

For full details see the example program
\ttt{example/fastjet\_example.cc} that is distributed with the
\ttt{FastJet} code. For the reader who is familiar with \ttt{KtJet}
\cite{KtJet}, the example program is also given as it would be written
for \ttt{KtJet}, in the file \ttt{example/ktjet\_example.cc}.

A simplified version of the \fastjet example program is given below.

\begin{verbatim}
#include "FjPseudoJet.hh"
#include "FjClusterSequence.hh"
using namespace std;

int main (int argc, char ** argv) {
  vector<FjPseudoJet> input_particles;
  
  // read in input particles
  double px, py , pz, E;
  while (cin >> px >> py >> pz >> E) {
    // pushd FjPseudoJet of (px,py,pz,E) onto back of the input_particles
    input_particles.push_back(FjPseudoJet(px,py,pz,E)); 
  }

  // run the jet clustering with option R=1.0 and strategy=Best
  double Rparam = 1.0;
  FjClusterSequence clust_seq(input_particles, Rparam, Best);

  // extract the inclusive jets with pt > 5 GeV, sorted by pt
  double ptmin = 5.0;
  vector<FjPseudoJet> inclusive_jets = clust_seq.inclusive_jets(ptmin);

  // extract the exclusive jets with dcut = 25 GeV^2 and sort them
  // in order of increasing pt
  double dcut = 25.0;
  vector<FjPseudoJet> exclusive_jets = sorted_by_pt(
                                        clust_seq.exclusive_jets(dcut));

  // print out the details for each jet
  for (unsigned int i = 0; i < exclusive_jets.size(); i++) {
    // get constituents of the jet. NB this is through a member function 
    // of clust_seq because that is where the information is held.
    vector<FjPseudoJet> constituents = clust_seq.constituents(exclusive_jets[i])

    printf("%5u %15.8f %15.8f %15.8f %8u\n",
           i, exclusive_jets[i].rap(), exclusive_jets[i].phi(),
           exclusive_jets[i].perp(), constituents.size());
  }
}
\end{verbatim}

%======================================================================
\section{Compilation notes}

In order to access the \ttt{NlnN} strategy the library needs to be
compiled with the Computational Geometry Algorithms Library \ttt{CGAL}
\cite{CGAL}. 

The user should first download \ttt{CGAL} from
\url{http://www.cgal.org/} and then compile and install it. Under
linux, she or he will have to ensure that an environment variable
\ttt{CGAL\_MAKEFILE} points to the Makefile generated by \ttt{CGAL} at
install time, containing various definitions of locations of include
files. \ttt{CGAL} will encourage the user to set this variable during
the install process.

The choice of whether or not CGAL will be used is made in the
Makefile in the top directory of the distribution, through the value
of the variable \ttt{USE\_CGAL} which should be either \ttt{yes} or
\ttt{no}. 

For more details see the \ttt{INSTALL} file in the top directory.




%======================================================================
\begin{thebibliography}{9}

\bibitem{fastjet}
  M.~Cacciari and G.~P.~Salam,
  %``Dispelling the N**3 myth for the k(t) jet-finder,''
  arXiv:hep-ph/0512210.
  %%CITATION = HEP-PH 0512210;%%
  %%Cited 1 time in SPIRES-HEP


\bibitem{ktexcl}
  S.~Catani, Y.~L.~Dokshitzer, M.~H.~Seymour and B.~R.~Webber,
  %``Longitudinally invariant K(t) clustering algorithms for hadron hadron
  %collisions,''
  Nucl.\ Phys.\ B {\bf 406}  (1993)  187.

\bibitem{ktincl}
  S.~D.~Ellis and D.~E.~Soper,
  %``Successive Combination Jet Algorithm For Hadron Collisions,''
  Phys.\ Rev.\ D {\bf 48} (1993) 3160 
  %\prd{48}{1993}{3160} 
  [hep-ph/9305266]. 
  %%CITATION = HEP-PH 9305266;%%
  



\bibitem{HepJet} \texttt{HepJet} is being planned in collaboration
  with authors of \ttt{KtJet} \cite{KtJet}.

%\cite{Blazey:2000qt}
\bibitem{RunII-jet-physics}
G.~C.~Blazey {\it et al.},
%``Run II jet physics,''
hep-ex/0005012.
%%CITATION = HEP-EX 0005012;%%

\bibitem{KtJet}
  \url{http://hepforge.cedar.ac.uk/ktjet/};
  J.~M.~Butterworth, J.~P.~Couchman, B.~E.~Cox and B.~M.~Waugh,
  %``KtJet: A C++ implementation of the K(T) clustering algorithm,''
  Comput.\ Phys.\ Commun.\  {\bf 153}, 85 (2003)
  [hep-ph/0210022].
  %%CITATION = HEP-PH 0210022;%%


\bibitem{CGAL}
%Andreas Fabri, Geert-Jan Giezeman, Lutz Kettner, Stefan Schirra and Sven Schonherr,
A.~Fabri {\it et al.},
%On the design of CGAL a computational geometry algorithms library
Softw.~Pract.~Exper.~ {\bf 30} (2000) 1167;
%Jean-Daniel Boissonnat, Olivier Devillers, Sylvain Pion, Monique Teillaud and Mariette Yvinec
J.-D.~Boissonnat {\it et al.},
% Triangulations in CGAL
Comp.~Geom.~{\bf 22} (2001) 5; \url{http://www.cgal.org/}

\end{thebibliography}


\end{document}
